% The appendix of the notes: Bilaga A, the method, written for the
% students.  This deck has no backed-claim chapters of its own: every
% claim it makes that cannot be checked at the keyboard is already
% answered in one of the lecture decks, and the two that are not are
% listed in Fortsatt arbete.  No \chapterprecis here -- that sentence
% belongs to the backed-claim chapters.

\chapter{Att kontrollera ett påstående}
\label{app:verifiera}

Övningen påstår, som allt undervisningsmaterial, en del.  Det mesta går
att pröva vid datorn: att en metod vägrar flytta bokmärket förbi sista
sidan visar övningen med en körning, och varifrån
\mintinline{python}{print} fick sitt svar innan klassen hade någon
\mintinline{python}{__str__} svarar klassen själv på
(\cref{sec:arv}).  Båda kan du göra om själv.  Men
några påståenden handlar om vad forskningen säger, och dem kan ingen
körning avgöra:
\begin{enumerate}
  \item att det lönar sig att göra ett eget försök innan du läser
    lösningsförslaget (inledningen);
  \item att nybörjare systematiskt missförstår vad en klass och ett
    objekt är, och att just de missförstånden är värda att rikta
    uppgifter mot (\cref{sec:tva-bocker} och
    \cref{sec:tva-exemplar});
  \item att inkapsling är en etablerad princip i objektorienterad
    programmering, och därför värd att öva (lärandemålen);
  \item att beskrivande namn hjälper den som läser koden
    (\cref{sec:granskning});
  \item att det är bra att inte upprepa sig i ett program --- principen
    \emph{DRY} --- och att den kommer från Hunt och Thomas
    (\cref{sec:granskning});
  \item att komposition är att föredra framför arv, och arv rätt val
    först när relationen verkligen är
    \foreignlanguage{english}{is-a} (\cref{sec:komposition});
  \item att en klass minskar risken för fel när flera delar av ett
    program ändrar samma data (\cref{sec:representation});
  \item att det utvecklar den egna designförmågan att granska kod som
    andra har skrivit (\cref{sec:granskning});
  \item att det är bättre att ett programfel märks direkt än att det
    passerar tyst (\cref{sec:sista-sidan}).
\end{enumerate}
Hur vet man om sådant stämmer?  Samma fråga möter dig varje gång du
läser en lärobok, en webbsida eller ett svar från en språkmodell ---
och varje gång du själv ska skriva en rapport.  Den här bilagan
sammanfattar arbetsgången, och \cref{tab:ovn-pastaenden} säger var
svaret på varje påstående finns.  Sex av de nio är redan prövade i
föreläsningarnas bilagor, och tabellen anger vilken; de tre sista är
inte prövade alls, och \cref{sec:ovn-fortsatt} säger vad som återstår.
Arbetsgången beskrivs utförligare i bilagan med samma namn i
föreläsningen \emph{Algoritmiskt tänkande}.

\begin{table}[!ht]
  \centering
  \caption{Övningens nio påståenden som frågor, och var svaren finns.
    Numren i första kolumnen är påståendenas nummer i bilagans
    inledande uppräkning.
    \enquote{Ja, med förbehåll} betyder att bilagan i fråga svarar ja
    och räknar upp de förbehåll som hör till svaret; de förbehållen
    avgör hur påståendet får användas.}
  \label{tab:ovn-pastaenden}
  \begin{tabular}{@{}lp{0.38\linewidth}p{0.32\linewidth}@{}}
    \toprule
    \# & Fråga & Svar \\
    \midrule
    1 & Leder det till bättre lärande att försöka lösa ett problem
        innan begreppet lärs ut? & Ja, med förbehåll; \emph{Upprepningar},
        bilaga~C \\
    2 & Missförstår nybörjare vad en klass och ett objekt är, och i så
        fall hur? & Ja, med förbehåll; \emph{Klasser och objekt},
        bilaga~B \\
    3 & Är inkapsling en etablerad princip i objektorienterad
        programmering? & Ja, med förbehåll; \emph{Klasser och objekt},
        bilaga~C \\
    4 & Hjälper beskrivande namn den som läser koden? & Delvis, och
        svagare än man tror, men namnen har stöd; \emph{Variabler och
        utskrifter}, bilaga~B \\
    5 & Är det bra att inte upprepa sig, och kommer DRY från Hunt och
        Thomas? & Ja, med förbehåll; \emph{Algoritmiskt tänkande},
        bilaga~E \\
    6 & Bör man föredra komposition framför arv? & Ja, med förbehåll;
        \emph{Praktiska tillämpningar av klasser}, bilaga~B \\
    7 & Minskar en klass risken för fel när flera delar av ett program
        ändrar samma data? & Oprövat; se \cref{sec:ovn-fortsatt} \\
    8 & Utvecklar det designförmågan att granska andras kod? &
        Oprövat; se \cref{sec:ovn-fortsatt} \\
    9 & Är det bättre att ett programfel märks direkt än att det
        passerar tyst? & Oprövat; se \cref{sec:ovn-fortsatt} \\
    \bottomrule
  \end{tabular}
\end{table}

\section{Gör påståendet till en fråga}
\label{sec:ovn-fraga}

Ett påstående går att kontrollera först när det är tydligt vad som
skulle kunna visa att det är \emph{fel}.  Formulera det därför som en
fråga med ett svar.  \enquote{Missförstår nybörjare vad en klass är?}
kan besvaras med ja eller nej; \enquote{klasser är bra} kan det inte,
eftersom \enquote{bra} är ett omdöme och inte ett faktum.  Ofta döljer
sig flera påståenden i ett: \enquote{DRY är en etablerad princip från
Hunt och Thomas} är två frågor --- kommer den från dem, och är den
etablerad --- och de kräver olika slags källor.

Lägg märke till vilket slags påståenden som hamnar i en sådan lista.
Det är inte
de som handlar om vad Python \emph{gör} --- att ett attribut med dubbla
understreck inte går att nå utifrån är en sak du kan köra --- utan de som
handlar om människor: vad nybörjare missförstår, vad som hjälper en
läsare, vad som lönar sig att göra först.  Varje gång ett
undervisningsmaterial säger \enquote{det här är bättre} har det gjort ett
sådant påstående.

\section{Ett svar som redan finns söks inte om}
\label{sec:ovn-atervinning}

Sex av de nio frågorna har redan ett svar, sökt en gång i den
föreläsning där påståendet först gjordes.  Då söks de inte om här.  I
stället citeras samma källa, anteckningen om hur den hittades och
lästes följer med till den här övningens litteraturlista, och texten
pekar på bilagan där svaret står --- med en mening om vad den bilagan
faktiskt besvarar, så att du vet vad du får om du går dit.

Det är inte en genväg utan hur ett kunskapsläge ska användas: samma
påstående prövas en gång, och den prövningen får bära överallt där
påståendet dyker upp.  Men det ställer ett krav: påståendet får inte
växa på vägen.  Bilagan om PEP~8 säger att \emph{namnen} har stöd och
att flera andra stilregler inte har det, och därför säger
\cref{sec:granskning} precis det och ingenting mer.

\section{Läs i fulltext, och skriv ner hur}
\label{sec:ovn-fulltext}

En träff i en databas är inte ett belägg.  Sammanfattningen
(\foreignlanguage{english}{abstract}) räcker inte heller: den säger vad
författarna vill ha sagt, inte vad de visat.  Läs hela texten och leta
upp \emph{det ställe} som styrker påståendet --- ett ordagrant citat med
sidnummer.  Hittar du inget sådant ställe stöder källan inte påståendet,
hur passande titeln än är.

Skriv sedan ner hur du gick till väga, så att någon annan kan följa
spåret.  I det här materialet står den anteckningen som en kommentar
ovanför varje post i litteraturlistans källfil, med fasta fält:
\begin{description}
  \item[\texttt{CLAIM}] vilket påstående källan ska styrka;
  \item[\texttt{FOUND-VIA}] hur den hittades --- sökfråga och databas,
    eller \enquote{känt dokument};
  \item[\texttt{PICKED}] varför just den, framför andra träffar;
  \item[\texttt{QUOTE}] det ordagranna citatet, med sida;
  \item[\texttt{VERIFIED}] att fulltexten lästs, och var;
  \item[\texttt{COUNTER}] motsökningen och vad den gav;
  \item[\texttt{DATE}] när.
\end{description}
Posterna i den här övningens litteraturlista är kopierade från
föreläsningarnas, och varje sådan kopia har fått två egna rader: en som
säger varifrån blocket är kopierat och vilken bilaga som bär prövningen,
och en som säger vad källan används till här och att den inte lästs om.
Föreläsningens egna anteckningar står kvar, men märkta med den
föreläsningens namn i stället för \enquote{här}, så att det går att se
vem som har gjort vad.  Utan de raderna hade det sett ut som om övningen
själv hade läst allt, och det hade varit osant.

\section{Fortsatt arbete}
\label{sec:ovn-fortsatt}

Tre av övningens påståenden är oprövade, och alla tre är av samma
slag: de låter som självklarheter och är det inte.

Det första är påståendet i \cref{sec:representation}: att en klass
minskar risken för fel när flera delar av ett program ändrar samma data.
Argumentet i lösningsförslaget är logiskt --- regeln bor på ett ställe i
stället för på många --- men påståendet att det \emph{i praktiken} ger
färre fel är empiriskt, och det är inte sökt.  Föreläsningen
\emph{Klasser och objekt} gör samma påstående, och inte heller där finns
det ett underlag; frågan att ställa är om inkapsling mätbart minskar
antalet defekter, och en motsökning ska leta efter studier som inte
finner någon skillnad.  Att bilagan om inkapsling
(\emph{Klasser och objekt}, bilaga~C) redan visar att principen följs
ofullständigt i praktiken är ett skäl att tro att svaret blir
\enquote{ja, med förbehåll} och inte \enquote{ja}.

Det andra är uppgiftsformen i \cref{sec:granskning}: att det utvecklar
den egna designförmågan att granska kod som andra har skrivit.  Här
finns ett helt forskningsfält --- kamratgranskning och kodgranskning i
programmeringsundervisning --- och tills det är genomsökt påstår texten
ingenting om nyttan, utan ber bara om granskningen.  Frågan att ställa
är om granskning av andras kod förbättrar den egna koden, och
motsökningen ska leta efter försök där effekten uteblev.

Det tredje är meningen i \cref{sec:sista-sidan} om att det är bättre
att ett programfel märks direkt än att det passerar tyst.  Rådet har
namn i litteraturen --- \foreignlanguage{english}{fail fast},
defensiv programmering --- och det finns en motsida: ett särfall som
avbryter ett program mitt i ett arbete kan lämna efter sig ett halvt
utfört arbete, och den som fångar allt för att slippa avbrottet är
tillbaka i det som föreläsningen \emph{Inmatning och felhantering},
bilaga~D, avråder från.  Frågan att ställa är om tidigt avbrott ger
färre eller billigare fel än att fortsätta med ett felaktigt värde.

Alla tre frågorna ska ställas ämnesdrivet, i samma databaser och med samma
motfrågor som de övriga bilagorna använder, och svaren ska skrivas ut
med sina förbehåll --- eller så ska påståendena strykas ur texten.
